\section{Background: fake news and simulation models for analysing its dissemination}

This section establishes the conceptual and methodological foundations of the study. It first explains why fake-news dissemination is mediated by source identity and social interaction, then introduces ABM and the ODD reporting protocol, and finally positions the proposed model relative to existing misinformation simulations.

\subsection{Fake news and source-mediated dissemination}

Fake news is not merely erroneous information; it is fabricated content presented in a news-like form and connected to incentives, attention, and institutional trust \citep{lazer2018,allcott2017}. Social platforms alter the route by which people encounter this content. Users select from sources, observe peers, and react to signals such as credibility, emotional tone, familiarity, and proximity. Homophilous communities can reinforce selective exposure \citep{delvicario2016}, and social influence can change evaluations independently of intrinsic quality \citep{muchnik2013}.

Source identity matters because users often lack the time or information needed to verify each claim. Aggregated judgments of source quality can support misinformation interventions \citep{pennycook2019}, whereas low-credibility content may receive coordinated amplification from automated accounts \citep{shao2018}. These mechanisms make source-level policies---visibility limits, recommendation rules, labels, and sanctions---plausible levers. They also create adaptation incentives: a malicious source may modify visible attributes to resemble a trusted source.

\subsection{Agent-based models as simulation models}

ABMs describe system outcomes as the emergent result of local decisions and interactions. They are useful when agents differ, remember past events, learn from contacts, and respond to interventions over time. In social-network research, \ifdefined\ANONYMOUS related prior applications are withheld for blinded review.\else ABMs have evaluated word-of-mouth (WOM) campaigns and participant selection \citep{leger2016,plaza2019}, message repetition on Twitter \citep{lopez2023}, and extensible architectures for social-network simulation \citep{leger2026}.\fi{} These applications illustrate how an executable mechanism can connect individual exposure and choice to population-level diffusion while retaining a controlled experimental design.

A credible ABM must make its entities, processes, scheduling, stochasticity, inputs, and outputs explicit. We follow the ODD protocol \citep{grimm2020}, which organises a model description into Overview, Design concepts, and Details. Figure~\ref{fig:odd} summarises how those components move from the scientific purpose and system boundary to behavioural assumptions and executable submodels. ODD supports inspection and replication, but does not itself establish that a model is valid for every use.

\begin{figure}[ht]
\centering
\includegraphics[width=.98\linewidth]{odd-protocol-abm.pdf}
\caption{Original synthesis of the ODD protocol for documenting an agent-based model. The three descriptive components support replication, software verification, and qualified validation.}
\label{fig:odd}
\end{figure}

Simulation evidence has a different status from observational or experimental evidence about people. A simulation supports conditional statements: given the encoded mechanisms and parameter ranges, a strategy produces a measured change. Verification asks whether the software implements that specification; validation asks whether the representation is sufficiently credible for its intended purpose. Neither process turns an exploratory model into a direct prediction of platform behaviour.

\subsection{Simulation models of misinformation}

Misinformation models range from epidemic analogies to network, opinion-dynamics, and cognitively richer agent systems. Table~\ref{tab:related-models} compares representative models by their behavioural focus, empirical input, intervention and validation. Competing-rumour models represent misinformation and fact-checking as diffusion processes \citep{sulis2020}; game-theoretic models examine deception and cooperation \citep{hills2018}; curation models connect belief, recommendation objectives and empirically observed Twitter data \citep{gausen2022}; and bot models study how network position changes indirect influence \citep{keijzer2021}. Recent work also reconstructs language and network histories through digital cloning \citep{puri2024}, examines moderation across platforms \citep{murdock2024}, and evaluates governance in decentralised systems \citep{malgonde2025}.

\begin{table*}[ht]
\centering
\caption{Positioning relative to representative misinformation simulation models. ``Empirical'' identifies inputs or validation data reported by each study; it does not imply predictive validity.}
\label{tab:related-models}
\small
\begin{tabular}{p{2.5cm}p{2.8cm}p{2.6cm}p{2.5cm}p{3.4cm}}
\toprule
Study & Main mechanism & Network/input & Intervention & Distinctive analytical focus \\
\midrule
\citet{sulis2020} & Competing misinformation and debunking states & Configurable synthetic networks & Fact-checking probability & Diffusion and forgetting \\
\citet{hills2018} & Evolutionary deception and cooperation & Random interactions; empirical pattern comparison & Behavioural costs & Deception types and cooperation \\
\citet{gausen2022} & Beliefs and newsfeed curation & Parameterisation and case comparisons using Twitter data & Curation objectives & Algorithmic misinformation--polarisation trade-off \\
\citet{keijzer2021} & Cultural influence by bots and humans & Synthetic network structures & Bot activity and connectivity & Direct versus indirect bot influence \\
\citet{puri2024} & Language-sensitive digital clones & Observed network and message histories & Counterfactual content & Reproduction of a particular online system \\
This study & Repeated source choice from accumulated evidence & Survey-informed source perceptions; synthetic directed networks & Source imitation, reach and recommendation policy & Adaptive source presentation with objective false-publication probability held separate \\
\bottomrule
\end{tabular}
\end{table*}


Our model occupies a complementary niche. It does not reconstruct a particular X network or classify text. Instead, it makes a source---rather than a message or belief state---the adaptive experimental actor, keeps objective false-publication behaviour separate from perceived attributes, and evaluates selective and combined source imitation. This makes it suitable for controlled scenario analysis and for distinguishing source-level success from system-level false-news exposure. Its novelty is therefore the experimental treatment of source presentation and platform-controllable reach, not higher-fidelity reconstruction of X.
